\usepackage[utf8]{inputenc}
\usepackage[T1]{fontenc}
\usepackage{lmodern}
\usepackage[margin=1in]{geometry}
\usepackage{amsmath,amssymb,amsthm}
\usepackage{enumitem}
\usepackage{booktabs}
\usepackage{xcolor}
\usepackage{tikz}
\usepackage[colorlinks=true,bookmarks=false,linkcolor=blue!50!black,urlcolor=blue!50!black]{hyperref}

\theoremstyle{plain}
\newtheorem{theorem}{Theorem}
\newtheorem{lemma}[theorem]{Lemma}
\newtheorem{claim}[theorem]{Claim}
\newtheorem{corollary}[theorem]{Corollary}
\theoremstyle{definition}
\newtheorem{definition}[theorem]{Definition}
\newtheorem{example}[theorem]{Example}
\theoremstyle{remark}
\newtheorem*{remark}{Remark}

\newcommand{\N}{\mathbb{N}}
\newcommand{\Z}{\mathbb{Z}}
\newcommand{\R}{\mathbb{R}}
\newcommand{\Q}{\mathbb{Q}}
\newcommand{\iv}[1]{[#1]}                 % Iverson bracket
\newcommand{\powerset}{\mathcal{P}}
\DeclareMathOperator{\ord}{ord}
\DeclareMathOperator{\gcdop}{gcd}

\setlength{\parindent}{0pt}
\setlength{\parskip}{0.6em}
\newcommand{\lesson}[2]{\href{../lessons/#1}{#2}}   % link to a lesson image

% Let long inline lists like \{1,2,5,6\} break across lines after a comma.
\mathchardef\mycomma="613B
\mathcode`,="8000
{\catcode`,=\active \gdef,{\mycomma\penalty0\hspace{0pt}}}
\emergencystretch=2em
